\documentclass[11pt,a4paper]{article}

% --- Packages ---
\usepackage[utf8]{inputenc}
\usepackage[T1]{fontenc}
\usepackage[english]{babel}
\usepackage{amsmath,amssymb,amsfonts,amsthm}
\usepackage{geometry}
\geometry{margin=1in}
\usepackage{graphicx}
\usepackage{hyperref}
\usepackage{xcolor}
\usepackage{cite}
\usepackage{booktabs}
\usepackage{microtype}
\usepackage{physics}
\usepackage{tikz}
\usetikzlibrary{babel}

\hypersetup{
    pdftitle={Scale-Time Algebra and Conformal Calculus: A New Mathematical Formalism for Contracting Systems},
    pdfauthor={Muhammed Furkan Haydari},
    pdfsubject={Mathematical Physics, Differential Geometry, Nilpotent Algebra, Scalor Space},
    pdfkeywords={Scale-Time Algebra, Scalor, Nilpotent Contraction Generator, Hubble Identity},
    colorlinks=true,
    linkcolor=blue!70!black,
    citecolor=red!70!black,
    urlcolor=blue!80!black
}

% --- Theorem Environments ---
\newtheorem{theorem}{Theorem}[section]
\newtheorem{definition}{Definition}[section]
\newtheorem{axiom}{Axiom}[section]
\newtheorem{corollary}{Corollary}[theorem]
\newtheorem{lemma}[theorem]{Lemma}
\newtheorem{remark}{Remark}[section]

\title{\textbf{Scale-Time Algebra and Conformal Calculus: \\ A New Mathematical Formalism for Contracting Systems}}

\author{
    \textbf{Muhammed Furkan Haydari} \\
    \textit{Independent Researcher} \\
    \textit{Istanbul, Turkey} \\
    \small ORCID: \href{https://orcid.org/0009-0003-2854-3947}{0009-0003-2854-3947}
}

\date{\today}

\begin{document}

\maketitle

\begin{abstract}
    Classical differential geometry and coordinate formalisms are constructed upon the premise that reference grids and measurement standards are rigid (time-invariant) manifolds. In this work, we formulate \textbf{Scale-Time Algebra} and its associated differential calculus to describe physical systems in which the background spacetime remains strictly static and Euclidean, while local atomic rulers and bound states undergo global conformal contraction.

    At the foundation of this formalism lies a novel nilpotent algebraic generator, \textbf{$\kappa$ ($\kappa \neq 0, \kappa^2 = 0$)}, structurally congruent with Clifford dual numbers and Grassmann exterior algebras. We define the \textbf{Scalor} space, which unifies static spatial coordinates with local conformal scale factors into a single hyper-complex element. By constructing the Scale-Lorentz group $\mathcal{SL}(3,1)$ representations, we derive the \textbf{Scale-Covariant Derivative ($\mathbb{D}_t$)} as the natural generalization of standard differential operators. As a direct analytic consequence of this calculus, the apparent cosmological Hubble relation ($v = H_0 d$) emerges directly from the Leibniz product rule applied to a contracting coordinate standard, demonstrating how scale evolution mimics kinematic recession without requiring metric expansion of space.
\end{abstract}

\vspace{1em}
\noindent\textbf{Keywords:} Scale-Time Algebra, Scalor, Contraction Generator, Scale-Covariant Derivative, Hubble Identity, Nilpotent Algebra, Clifford Algebra.

\vspace{1em}
\hrule
\vspace{0.8em}
\tableofcontents
\vspace{0.8em}
\hrule
\vspace{1.5em}

\section{Introduction and Mathematical Motivation}
\label{sec:giris}

Standard differential geometry and metric calculus are founded upon the implicit assumption that coordinate meshes and measuring rods are time-independent (rigid). The real line $\mathbb{R}$ establishes distance between two points as $d(x_1, x_2) = |x_1 - x_2|$, while the complex field $\mathbb{C} = \{x + i y \mid i^2 = -1\}$ parameterizes two-dimensional orthogonal rotations and oscillatory wave dynamics.

However, as emphasized by Christof Wetterich~\cite{wetterich2013} and Fred Hoyle~\cite{hoyle1964}, an exact \textbf{conformal equivalence} exists between an expanding FLRW metric ($g_{\mu\nu} = a^2(\eta) \eta_{\mu\nu}$) and a static Minkowski background where local atomic length standards undergo continuous contraction ($\lambda(t) \equiv a^{-1}(t)$). Whereas standard cosmology attributes this degree of freedom to metric tensor expansion, here we encode scale evolution directly as a nilpotent component of the underlying algebraic structure.

\subsection{Conformal Geometric Reduction of Fundamental Quantities}
\label{sec:geometrik_indirgeme}

To formulate the interaction between spacetime geometry and dynamical fields within a single unified framework, time ($t$), position ($\mathbf{x}$), mass ($M$), and the speed of light ($c$) are represented differentially on a quiescent background in terms of the contraction dynamics ($\lambda(t)$ and nilpotent generator $\kappa$):
\begin{enumerate}
    \item \textbf{Speed of Light ($c$):} Governs the transverse phase propagation velocity within the stationary vacuum continuum and the local quantum phase horizon ceiling ($v_c = 1/\sqrt{\varepsilon_0\mu_0} = c$).
    \item \textbf{Time ($t$):} Represents the differential logarithmic rhythm recorded by atomic quantum phase oscillations along the scale contraction trajectory ($dt \equiv -\frac{1}{H_0} d\ln\lambda$).
    \item \textbf{Position / Coordinate Distance ($R$):} The dimensionless ratio between the invariant coordinate distance in the static Euclidean continuum and the local contracting ruler ($R(t) = \mathbf{X}_0 / \lambda(t)$).
    \item \textbf{Mass ($M$):} The net Gauss flux of the convective vacuum acceleration field across a closed surface within the Painlevé--Gullstrand geometry ($M \equiv \frac{1}{4\pi G} \oint_{\partial \Omega} |(\mathbf{v}_{\mathrm{inflow}} \cdot \nabla) \mathbf{v}_{\mathrm{inflow}}| \, dA$).
\end{enumerate}
With these geometric identifications, the dynamics of contracting systems can be analyzed within a scale-covariant differential calculus.

\section{Fundamental Algebraic Structure: $\kappa$ Contraction Generator and Scalors}
\label{sec:cebir}

\begin{definition}[$\kappa$ Contraction Unit and Nilpotent Dual Algebra]
    \label{def:kappa_nilpotent_sc}
    Let $\kappa$ denote an algebraic unit outside the real field $\mathbb{R}$, satisfying the nilpotent relation:
    \begin{equation}
        \kappa \neq 0, \quad \kappa^2 = 0.
        \label{eq:kappa_tanim}
    \end{equation}
    Historically formulated by William Kingdon Clifford (1873) as the ring of dual numbers ($\mathbb{D} = \{a + b\kappa \mid a, b \in \mathbb{R}, \kappa^2 = 0\}$) and developed by Eduard Study (1903) for kinematic geometry, this structure is physically extended in Scale-Time Calculus to model a 5-dimensional conformal super-manifold $\mathcal{M}^{(4|1)}$, attaching an internal scale contraction fiber to 4-dimensional static spacetime. $\kappa$ acts as the generator of internal conformal scale contraction, preserving the static spatial coordinate while encoding the radial contraction gradient.
\end{definition}

Just as the imaginary unit $i$ endows real quantities with rotational and phase dimensions:
$$z = x + i y \in \mathbb{C}, \quad e^{i\theta} = \cos\theta + i \sin\theta,$$
the unit $\kappa$ equips a quantity with an intrinsic scale dimension.

\begin{definition}[Scalor]
    \label{def:skalor_sc}
    A \textbf{Scalor} is a hyper-complex element combining an invariant spatial position vector $\mathbf{x} \in \mathbb{R}^3$ with a time-dependent atomic scale factor $\lambda(t) \in \mathbb{R}^+$ via the unit $\kappa$:
    \begin{equation}
        \mathbb{S}(t) \equiv \mathbf{x} + \kappa \cdot \lambda(t),
        \label{eq:skalor_tanim}
    \end{equation}
    where $\mathbf{x}$ denotes the coordinate position in the invariant static background, and $\lambda(t) = e^{-H_0 t}$ represents the internal atomic length standard.
\end{definition}

\subsection{Operations in Scalor Algebra}
For two Scalors $\mathbb{S}_1 = \mathbf{x}_1 + \kappa \lambda_1$ and $\mathbb{S}_2 = \mathbf{x}_2 + \kappa \lambda_2$:
\begin{enumerate}
    \item \textbf{Addition:}
          \begin{equation}
              \mathbb{S}_1 + \mathbb{S}_2 = (\mathbf{x}_1 + \mathbf{x}_2) + \kappa (\lambda_1 + \lambda_2).
          \end{equation}
    \item \textbf{Scale-Modulated Product:}
          Due to nilpotency $\kappa^2 = 0$:
          \begin{equation}
              \mathbb{S}_1 \odot \mathbb{S}_2 = (\mathbf{x}_1 \cdot \mathbf{x}_2) + \kappa (\lambda_1 \mathbf{x}_2 + \lambda_2 \mathbf{x}_1).
          \end{equation}
    \item \textbf{Conformal Exponential Expansion:}
          \begin{equation}
              e^{\kappa \theta} = 1 + \kappa \theta.
              \label{eq:kappa_ustel}
          \end{equation}
\end{enumerate}

\subsection{Nilpotent Algebra ($\kappa^2 = 0$), Dual Numbers, and Annihilation}
\label{sec:nilpotent_annihilation}

In standard real analysis, the mapping $f(x) = x^2$ is two-to-one; for instance $(\pm 2)^2 = 4$, and the only real number whose square vanishes is zero ($x^2 = 0 \iff x = 0$).

In contrast, within Clifford dual numbers (introduced by W.~K.~Clifford in 1873) and Hermann Grassmann's exterior algebra, non-zero hyper-complex units exist whose square is identically zero. The unit $\kappa$ belongs strictly to this family:
\begin{equation}
    \kappa \neq 0, \quad \text{yet} \quad \kappa^2 = 0.
\end{equation}
For a general Scalor $z = a + b\kappa$ with $a, b \in \mathbb{R}$:
\begin{equation}
    z^2 = (a + b\kappa)^2 = a^2 + 2ab\kappa + b^2 \kappa^2 = a^2 + 2ab\kappa.
\end{equation}
For pure nilpotent states ($a = 0$):
\begin{equation}
    z = b\kappa \implies z^2 = b^2 \kappa^2 \equiv 0 \quad (\forall b \in \mathbb{R}, \; b \neq 0).
\end{equation}
This establishes the existence of infinitely many non-zero elements whose square vanishes identically.

\subsubsection{Physical Manifestation: Nilpotent Ideals and Chiral Representation}
In quantum field theory, fermionic fields are parameterized via nilpotent Grassmann variables ($\theta_i \theta_j = -\theta_j \theta_i \implies \theta^2 = 0$), providing the algebraic foundation for the Pauli Exclusion Principle and Berezin integration~\cite{berezin1966}.

In the Scale-Time formalism, representing chiral vacuum excitations via the nilpotent generator yields:
\begin{equation}
    \boldsymbol{\omega}_+ = +\boldsymbol{\omega}_0 \kappa, \quad \boldsymbol{\omega}_- = -\boldsymbol{\omega}_0 \kappa.
\end{equation}
The superposition of these opposing chiral states produces a nilpotent contact interaction:
\begin{equation}
    \mathbb{S}_{\mathrm{interaction}} \propto (\kappa) \cdot (\kappa) = \kappa^2 \equiv 0.
\end{equation}
Algebraically, $\kappa^2 = 0$ guarantees that second-order differential forms in Scalor space are automatically closed, and chiral projection operators $\mathcal{P}_\pm \equiv \frac{1}{2}(1 \pm \kappa)$ form a nilpotent ideal. This ensures scale conservation without requiring expanding spatial manifolds.

\subsection{Group-Theoretic Structure: The Scale-Lorentz Group $\mathcal{SL}(3,1)$}
In standard relativistic physics, spacetime symmetries are governed by the Poincaré group $\mathcal{P} = SO(3,1) \ltimes \mathbb{R}^{3,1}$, which presupposes rigid measurement standards. In Scale-Time Algebra, this structure is extended by the nilpotent generator $\kappa$ into the \textbf{Scale-Lorentz Group} $\mathcal{SL}(3,1) \equiv SO(3,1) \ltimes \mathbb{R}^+_\kappa$.

The commutation relations between Dirac matrices $\gamma^\mu$ and $\kappa$ in the Clifford algebra read:
\begin{equation}
    \{\gamma^\mu, \gamma^\nu\} = 2\eta^{\mu\nu}, \quad [\kappa, \gamma^\mu] = -\kappa \gamma^\mu, \quad \kappa^2 = 0.
    \label{eq:clifford_scale}
\end{equation}
This algebra ensures that quantum spinors preserve their phase angles under contraction while their radial norms scale with $\lambda(t)$. The Scale-Covariant Derivative $\mathbb{D}_\mu \equiv \partial_\mu + \kappa A_\mu^{\mathrm{scale}}$ acts as a gauge-covariant derivative for scale transformations, where $A_0^{\mathrm{scale}} = H_0/c$ is the universal scale gauge potential.

\section{Ontological Derivation of Time and Position from Scalor Algebra}
\label{sec:zaman_konum_turetilis}

In Newtonian physics and General Relativity, space $\mathbf{x}$ and time $t$ are posited as primitive coordinate backgrounds. Motivated by the "Problem of Time" in canonical quantum gravity (the absence of time in the Wheeler--DeWitt equation) and singularity paradoxes, space and time should be viewed as derived, secondary quantities.

In Scale-Time Algebra, the universe possesses a single invariant: \textbf{The scaleless, frictionless, static continuum $\mathcal{E}_\infty$}. A physical observer or bound matter state is a Scalor defined on this continuum by $\mathbb{S} = \mathbf{X}_0 + \kappa \lambda$. From this foundation, time and position are derived rigorously:

\begin{axiom}[Static Continuum and Nilpotent Existence]
    \label{ax:statik_sureklilik_sc}
    The universal background is an absolute, static, and metric-free Euclidean continuum ($\mathbf{X}_0 \in \mathcal{E}_\infty$, $\dot{\mathbf{X}}_0 = 0$). Matter exists as a bound state on this continuum carrying a nilpotent contraction charge $\kappa^2 = 0$.
\end{axiom}

\begin{theorem}[Derivation of Time: Time as the Contraction Phase Log]
    \label{thm:zaman_turetilis}
    Physical time $t$ is derived as the integral of the logarithmic contraction rate of the internal nilpotent scale $\lambda$:
    \begin{equation}
        dt \equiv -\frac{1}{H_0} \frac{d\lambda}{\lambda} = -\frac{1}{H_0} d(\ln \lambda) \implies t(\lambda) = -\frac{1}{H_0} \ln\left(\frac{\lambda}{\lambda_0}\right).
        \label{eq:zaman_turetilis}
    \end{equation}
    This directly determines the scale evolution of matter:
    \begin{equation}
        \lambda(t) = \lambda_0 e^{-H_0 t}.
    \end{equation}
\end{theorem}

\begin{proof}
    A physical observer cannot measure the "flow of time" directly; an observer only registers the internal quantum frequency of bound systems (atoms, clocks, orbital periods): $\omega(\lambda) = \frac{\hbar}{m(\lambda) \lambda^2}$. In any contracting system, the atomic frequency accelerates as $\omega(t) = \omega_0 e^{H_0 t}$. Establishing a bijection between the cumulative phase count $\Phi = \int \omega(t) dt$ and scale contraction $d\lambda/\lambda$, we find $dt \propto -d\lambda/\lambda$. Normalizing by the universal vacuum flow constant $H_0^{-1}$ yields Theorem~\ref{thm:zaman_turetilis}.
\end{proof}

\begin{corollary}[Geometric Origin of the Arrow of Time]
    \label{cor:zaman_oku_sc}
    The arrow of time originates from the strict monotonicity of atomic scale contraction ($\dot{\lambda} < 0$). The monotonic evolution $\lambda(t) = \lambda_0 e^{-H_0 t}$ establishes a smooth, non-singular affine parameterization ($t \in \mathbb{R}$) on the manifold. In a rigid reference frame where contraction is absent ($\dot{\lambda} = 0$), the differential scale variation vanishes ($dt = 0$), and local clocks experience no frequency drift.
\end{corollary}

\begin{theorem}[Derivation of Position and Distance: Distance as Scale Ratio]
    \label{thm:konum_turetilis}
    Observed physical position and distance $\mathbf{R}_{\mathrm{obs}}$ correspond to the algebraic projection of the static background point $\mathbf{X}_0$ onto the observer's contracted ruler scale $\lambda(t)$:
    \begin{equation}
        \mathbf{R}_{\mathrm{obs}}(t) \equiv \frac{\mathbf{X}_0}{\lambda(t)} = \frac{\mathbf{X}_0}{\lambda_0} e^{H_0 t}.
        \label{eq:konum_turetilis}
    \end{equation}
\end{theorem}

\begin{proof}
    An observer measures space not via abstract meters, but in units of their own atomic radius $r_{\mathrm{atom}}(t) = r_0 \lambda(t)$. The number of atomic ruler units fitting between two stationary points $(\mathbf{X}_1, \mathbf{X}_2)$ in the static background is:
    \begin{equation}
        N(t) \equiv \frac{|\mathbf{X}_2 - \mathbf{X}_1|}{r_{\mathrm{atom}}(t)} = \frac{|\mathbf{X}_0|}{r_0 \lambda(t)} = \frac{|\mathbf{X}_0|}{r_0 \lambda_0} e^{H_0 t} = N_0 e^{H_0 t}.
        \label{eq:cetvel_sayimi_ispat_sc}
    \end{equation}
    Although the absolute background distance $|\mathbf{X}_0|$ remains invariant ($\dot{\mathbf{X}}_0 = 0$), as the atomic rulers contract ($r_{\mathrm{atom}} \to 0$), the count of atomic units spanning the interval grows exponentially. An observer unaware of their own shrinking interprets this increasing ruler count as an apparent spatial recession ($\mathbf{R}_{\mathrm{obs}}$ growing) and recession velocity ($v = H_0 d$).
\end{proof}

\begin{quote}
    \small
    \textbf{Epistemological Synthesis: Resolution of the Newton vs. Leibniz--Mach Debate:} \\
    While Newton regarded space as an absolute container, Leibniz and Mach argued that space is merely the relational configuration among bodies. Scale-Time Calculus unifies this historical debate:
    \begin{itemize}
        \item \textbf{Background (Newtonian Absolute):} Space itself ($\mathcal{E}_\infty$) is an absolute, quiescent, uncurved, and non-expanding Euclidean continuum.
        \item \textbf{Distance and Separation (Machian Relational):} All perceived and measured distances are relational projections of a contracting atomic network ($R \sim 1/\lambda$).
    \end{itemize}
    This ontological duality clarifies two limiting regimes: In a local laboratory, matter contraction cannot be detected by internal comparisons (since both the sample and the ruler contract concurrently):
    \begin{equation}
        \lim_{\lambda(t) \to 0} \frac{L(t)}{r_{\mathrm{ruler}}(t)} = \lim_{\lambda \to 0} \frac{L_0 \lambda(t)}{1\text{ m} \cdot \lambda(t)} = L_0 = \text{Constant},
        \label{eq:lhopital_oran_sc}
    \end{equation}
    preserving exact internal scale invariance.

    Conversely, the conventional observable universe horizon ($\sim 93$ billion light-years) is not a physical boundary of space nor the total cosmic volume. The "light-year" is a relative ruler locked to the contraction cadence of local atomic clocks ($1\text{ ly} \equiv c \times 1\text{ local atomic year}$). Because the background continuum ($\mathcal{E}_\infty$) is strictly static and infinite ($X = \text{const}$), measuring an invariant Euclidean separation with shrinking atomic rulers yields:
    \begin{equation}
        d_{\mathrm{obs}}(t) = \frac{X}{\lambda(t)} \xrightarrow{\lambda(t) \to 0} \infty.
        \label{eq:sonsuz_arka_plan_iraksama_sc}
    \end{equation}
    The universe does not expand; rather, shrinking rulers register the quiescent infinity as an apparent expansion. The 93-billion-light-year boundary is merely the optical horizon integrated backward over our contracted coordinate window; beyond it, absolute static Euclidean space extends indefinitely.
\end{quote}

\begin{theorem}[Planck Length and Planck Time: The Vacuum Saturation Threshold]
    \label{thm:planck_ontolojisi}
    In conventional physics, the Planck length ($\ell_P \equiv \sqrt{\hbar G/c^3} \approx 1.6 \times 10^{-35}\text{ m}$) and Planck time ($t_P \equiv \ell_P/c \approx 5.4 \times 10^{-44}\text{ s}$) are often viewed as discreteness thresholds or spacetime foam. In Scale-Time Calculus, these quantities are re-identified ontologically:
    \begin{enumerate}
        \item \textbf{Planck Length ($\ell_P$):} The background vacuum $\mathcal{E}_\infty$ is strictly continuous and smooth; space has no discrete pixels. Rather, $\ell_P$ is the \textbf{Topological Saturation Threshold} where the Compton wavelength ($\lambda_C = \hbar/mc$) of a matter vortex intersects its vacuum flow event horizon ($r_s = 2Gm/c^2$):
              \begin{equation}
                  \lambda_C \sim r_s \implies \frac{\hbar}{mc} \sim \frac{Gm}{c^2} \implies r_{\mathrm{min}} \sim \sqrt{\frac{\hbar G}{c^3}} \equiv \ell_P.
                  \label{eq:planck_uzunluk_sc}
              \end{equation}
              Because the maximum tension force the vacuum can sustain is bounded by $F_{\mathrm{Planck}} = c^4/G \approx 1.21 \times 10^{44}\text{ N}$, no physical system can be compressed below the radius $\ell_P$.
        \item \textbf{Planck Time ($t_P$):} By the contraction log ($dt \equiv -\frac{1}{H_0} d\ln\lambda$), $t_P = \ell_P/c$ represents the \textbf{period of a single complete phase revolution of the most saturated physical quantum oscillator possible}:
              \begin{equation}
                  t_P = \frac{\ell_P}{c} = \sqrt{\frac{\hbar G}{c^5}} \implies \omega_{\mathrm{max}} = \frac{1}{t_P} \approx 1.85 \times 10^{43}\text{ s}^{-1}.
                  \label{eq:planck_zaman_sc}
              \end{equation}
              Intervals shorter than $t_P$ are physically meaningless because no physical clock mechanism in nature can oscillate at a higher frequency.
    \end{enumerate}
\end{theorem}

\begin{corollary}[Dirac Large Numbers and the Scale Contraction Factor]
    \label{cor:dirac_buyuk_sayilar_sc}
    The large ratio between the Bohr radius and the Planck length ($a_0 / \ell_P \sim 10^{25}$) represents the \textbf{cumulative logarithmic phase distance traversed by the universal contraction engine ($e^{-H_0 t}$) relative to the vacuum saturation ceiling $\ell_P$}.
\end{corollary}

\begin{corollary}[Timescale to Planck Saturation: $\sim 785$ Billion Years]
    \label{cor:planck_boyuna_inme_zamani_sc}
    As atomic radii contract according to $r_{\mathrm{atom}}(t) = a_0 e^{-H_0 \Delta t}$, the time required to reach the mechanical saturation limit $\ell_P$ is:
    \begin{equation}
        \Delta t_{\mathrm{Planck}} = \frac{1}{H_0} \ln\left( \frac{a_0}{\ell_P} \right) = t_H \ln\left( \frac{a_0}{\ell_P} \right).
        \label{eq:planck_zamani_hesap_sc}
    \end{equation}
    Taking $t_H \equiv 1/H_0 \approx 13.90\text{ Gyr}$, $a_0 \approx 5.29 \times 10^{-11}\text{ m}$, and $\ell_P \approx 1.62 \times 10^{-35}\text{ m}$:
    \begin{equation}
        \frac{a_0}{\ell_P} \approx 3.27 \times 10^{24} \implies \ln\left( \frac{a_0}{\ell_P} \right) \approx 56.45,
    \end{equation}
    yielding a future duration of:
    \begin{equation}
        \Delta t_{\mathrm{Planck}} \approx 56.45 \times 13.90\text{ Gyr} \approx \mathbf{785\text{ Billion Years}},
        \label{eq:785_milyar_yil_sc}
    \end{equation}
    (for a proton radius $r_p \approx 0.84\text{ fm}$, $\Delta t \approx 45.4 \, t_H \approx \mathbf{631\text{ Billion Years}}$). This confirms that 3D atomic physics will operate uninterrupted for roughly $57$ times the current age of the universe, after which matter transitions into a saturated 1D flux regime at the tension limit ($c^4/G$).
\end{corollary}

\begin{theorem}[Derivation of Mass: Mass as Vacuum Absorption Throughput]
    \label{thm:kutle_turetilis}
    Mass $M$ is defined as the \textbf{vacuum flux rate} drawn inward per unit time from the surrounding static vacuum continuum by contracting matter:
    \begin{equation}
        M \equiv \frac{1}{4\pi G} \oint_{\partial \Omega} v_{\mathrm{inflow}}^2 \, d\mathbf{A},
        \label{eq:kutle_tanim}
    \end{equation}
    where $v_{\mathrm{inflow}}(r) = -\sqrt{2GM/r}$ is the radial velocity field. Mass-energy equivalence ($E = mc^2$) reflects the total kinetic energy of this inward vacuum flow collapsing at the characteristic radial contraction speed $v_c = c$.
\end{theorem}

\begin{corollary}[Invariance of Dimensionless Ratios and Laboratory Null Tests]
    \label{cor:boyutsuz_oranlar_sc}
    Any laboratory measurement of two masses yields the dimensionless ratio $\mathcal{R}_m \equiv m_1 / m_2$. Because all bound matter states scale synchronously according to $m(t) = m_0 e^{H_0 t}$:
    \begin{equation}
        \frac{m_1(t)}{m_2(t)} = \frac{m_{1,0} e^{H_0 t}}{m_{2,0} e^{H_0 t}} = \frac{m_{1,0}}{m_{2,0}} = \text{const} \implies \frac{d}{dt}\left(\frac{m_1}{m_2}\right) \equiv 0.
        \label{eq:boyutsuz_oran_sc}
    \end{equation}
    Likewise, $\dot{\alpha}/\alpha = 0$ and $\dot{\mu}/\mu = 0$ (where $\mu = m_p/m_e$). Because local dimensionless ratios remain invariant, Eötvös torsion balances and the MICROSCOPE satellite experiment observe differential acceleration consistent with zero to within $\eta < 10^{-15}$.
\end{corollary}

\section{Scale-Covariant Calculus}
\label{sec:kalkulus}

In classical mechanics, particle velocity $\mathbf{v} = \frac{d\mathbf{x}}{dt}$ is defined via the standard Newtonian derivative. When the observer's measurement units contract with time, this definition must be extended.

\begin{definition}[Scale-Covariant Derivative]
    \label{def:olcek_kovaryant_turev_sc}
    The \textbf{Scale-Covariant Derivative Operator} $\mathbb{D}_t$, unifying the static coordinate derivative with internal scale contraction, is defined by:
    \begin{equation}
        \mathbb{D}_t \equiv \frac{d}{dt} - \kappa \left( \frac{\dot{\lambda}(t)}{\lambda(t)} \right) = \frac{d}{dt} + \kappa H_0,
        \label{eq:olcek_turev_operatoru}
    \end{equation}
    where $H_0 \equiv -\dot{\lambda}/\lambda$ is the universal contraction rate.
\end{definition}

\subsection{Fundamental Theorems of Scale Calculus}

\begin{theorem}[Relative Velocity and the Hubble Identity]
    \label{thm:hubble_teoremi}
    Let $\mathbf{R}(t) \equiv \frac{\mathbf{x}}{\lambda(t)}$ denote the physical distance between two bound states at rest in the static background ($d\mathbf{x}/dt = 0$), measured with a contracting local ruler $\lambda(t)$. The time derivative of this distance is:
    \begin{equation}
        \mathbf{v}_{\mathrm{obs}} = \frac{d\mathbf{R}}{dt} = \frac{\mathbf{v}_{\mathrm{trans}}}{\lambda(t)} + H_0 \mathbf{R}(t),
        \label{eq:hubble_ispat}
    \end{equation}
    which, for zero translational coordinate velocity ($\mathbf{v}_{\mathrm{trans}} = 0$), reduces directly to:
    \begin{equation}
        \mathbf{v}_{\mathrm{obs}} = H_0 \mathbf{R}.
        \label{eq:hubble_saf}
    \end{equation}
\end{theorem}

\begin{proof}
    Applying the product rule of calculus to $\mathbf{R}(t) = \mathbf{x}(t) \cdot \lambda^{-1}(t)$:
    \begin{align*}
        \mathbf{v}_{\mathrm{obs}} & = \frac{d}{dt} \left( \mathbf{x}(t) \lambda^{-1}(t) \right)                                                                            \\
                                  & = \frac{d\mathbf{x}}{dt} \lambda^{-1}(t) + \mathbf{x}(t) \frac{d(\lambda^{-1}(t))}{dt}                                                 \\
                                  & = \frac{\mathbf{v}_{\mathrm{trans}}}{\lambda(t)} + \mathbf{x}(t) \left( -\frac{\dot{\lambda}(t)}{\lambda^2(t)} \right)                 \\
                                  & = \frac{\mathbf{v}_{\mathrm{trans}}}{\lambda(t)} + \left(-\frac{\dot{\lambda}(t)}{\lambda(t)}\right) \frac{\mathbf{x}(t)}{\lambda(t)}.
    \end{align*}
    Substituting $H_0 \equiv -\frac{\dot{\lambda}(t)}{\lambda(t)}$ and $\mathbf{R}(t) = \frac{\mathbf{x}(t)}{\lambda(t)}$:
    $$\mathbf{v}_{\mathrm{obs}} = \frac{\mathbf{v}_{\mathrm{trans}}}{\lambda(t)} + H_0 \mathbf{R}(t).$$
    For bodies at rest in coordinate space ($\mathbf{v}_{\mathrm{trans}} = 0$):
    $$\mathbf{v}_{\mathrm{obs}} = H_0 \mathbf{R}.$$
    This completes the proof.
\end{proof}

\section{Three Distinct Velocity Regimes}
\label{sec:uc_hiz}

In Scale Calculus, the concept of velocity separates into distinct categories:

\begin{table}[h!]
    \centering
    \small
    \resizebox{\textwidth}{!}{
    \begin{tabular}{@{}llll@{}}
        \toprule
        \textbf{Velocity Type}                 & \textbf{Symbol}               & \textbf{Mathematical Definition}                                                & \textbf{Physical Meaning}                         \\ \midrule
        \textbf{Coordinate Velocity}           & $\mathbf{v}_{\mathrm{trans}}$ & $\frac{d\mathbf{x}}{dt}$ ($v < c$)                                              & Matter translation on static background          \\
        \textbf{Intrinsic Phase Velocity}      & $v_c$                         & $1/\sqrt{\varepsilon_0\mu_0} = c$                                               & Elastic vacuum phase propagation ceiling          \\
        \textbf{Vacuum Phase (Light) Velocity} & $v_{\mathrm{phase}}$          & $\frac{\omega}{|\mathbf{k}|} = c$                                               & Phase speed relative to contracting observer      \\
        \textbf{Apparent (Scale) Velocity}     & $\mathbf{v}_{\mathrm{obs}}$   & $\mathbb{D}_t \mathbf{R} = \frac{v_{\mathrm{trans}}}{\lambda} + H_0 \mathbf{R}$ & Apparent velocity recorded by contracting ruler   \\ \bottomrule
    \end{tabular}
    }
    \caption{Hierarchy of velocity definitions in Scale-Time Algebra.}
    \label{tab:hiz_hiyerarsisi}
\end{table}

\subsection{Coordinate Separation Velocity (\texorpdfstring{$v_{\mathrm{sep}} \le 2c$}{vsep <= 2c})}
\label{sec:ayrilma_hizi}

The relativistic postulate that no physical signal can exceed $c$ applies strictly to \textbf{local information and particle propagation through the vacuum} ($v_{\mathrm{local}} \le c$). On our static Euclidean manifold ($\mathcal{E}_\infty$), the rate of separation between two independent wavefronts is an unconstrained kinematic quantity:

\begin{definition}[Coordinate Separation Velocity - $v_{\mathrm{sep}}$]
    \label{def:koordinat_ayrilma_hizi_sc}
    Let two photons propagate in opposite directions ($180^\circ$) across the static vacuum with positions $\mathbf{x}_1(t) = -c t \, \hat{n}$ and $\mathbf{x}_2(t) = +c t \, \hat{n}$. The Euclidean distance between them is $D(t) \equiv |\mathbf{x}_2(t) - \mathbf{x}_1(t)| = 2ct$. The coordinate separation velocity is:
    \begin{equation}
        v_{\mathrm{sep}} \equiv \frac{d D}{dt} = 2c.
        \label{eq:ayrilma_hizi_2c}
    \end{equation}
    This does not violate causality, as neither individual wavefront propagates through the vacuum at a speed exceeding $c$.
\end{definition}

\subsection{The Scale-Locking Principle}
\label{sec:olceksel_kilitlenme}

The stability of bound states across atomic and celestial mechanics is governed by the \textbf{Scale-Locking} condition:
\begin{equation}
    \mathbb{D}_t \left( \frac{\mathbf{x}(t)}{\lambda(t)} \right) = 0 \iff \mathbf{R}(t) = \text{const}.
    \label{eq:scale_locking_tanim}
\end{equation}
As a system contracts ($\dot{\lambda} < 0$), its local length standards shrink at the identical rate, maintaining a constant dimensionless coordinate $\mathbf{R}(t)$ in its comoving frame.

\section{Scale-Covariant Wave Equation}
\label{sec:dalga_denklemi}

Generalizing the d'Alembert operator via $\mathbb{D}_\mu = (\frac{1}{v_c}\mathbb{D}_t, \nabla)$:
\begin{equation}
    \mathbb{D}^\mu \mathbb{D}_\mu \Psi = \left( -\frac{1}{v_c^2} \left(\frac{\partial}{\partial t} + \kappa H_0\right)^2 + \nabla^2 \right) \Psi = 0.
    \label{eq:genel_dalga_operatoru}
\end{equation}
Applying the nilpotency relation $\kappa^2 = 0$, the wave operator simplifies to:
\begin{equation}
    \left( -\frac{1}{v_c^2} \frac{\partial^2}{\partial t^2} - \frac{2\kappa H_0}{v_c^2} \frac{\partial}{\partial t} + \nabla^2 \right) \Psi = 0.
    \label{eq:acik_dalga}
\end{equation}
The cross term $\frac{2\kappa H_0}{v_c^2} \frac{\partial}{\partial t}$ is the geometric scale-damping term that accounts for why photons propagating through a quiescent vacuum are registered as \textbf{redshifted} by contracting atomic detectors.

\section{Scale-Covariant Phase Space and Heisenberg Uncertainty}
\label{sec:faz_uzayi_heisenberg}

By Liouville's theorem in classical and quantum statistical mechanics, phase space volume is conserved under Hamiltonian time evolution: $d\Gamma = d^3\mathbf{x} \, d^3\mathbf{p} = \text{const}$. In a framework where the observer's spatial scale contracts ($\lambda(t) \to 0$):

\begin{theorem}[Scale-Invariance of Phase Space Volume]
    \label{thm:faz_uzayi_degismezligi_sc}
    As the characteristic spatial dimension of a bound quantum state contracts ($\Delta x(t) = \lambda(t) \Delta x_0$), its momentum uncertainty expands reciprocally via the de Broglie relation ($\lambda_{\mathrm{dB}} = \hbar / p$):
    \begin{equation}
        \Delta p(t) = \frac{\Delta p_0}{\lambda(t)}.
        \label{eq:momentum_genlesme}
    \end{equation}
    Consequently, the differential volume element of Scalor phase space is strictly \textbf{scale-invariant}:
    \begin{equation}
        d\Gamma(t) \equiv d^3\mathbf{x}(t) \, d^3\mathbf{p}(t) = \left( \lambda^3(t) d^3\mathbf{x}_0 \right) \cdot \left( \lambda^{-3}(t) d^3\mathbf{p}_0 \right) = d^3\mathbf{x}_0 \, d^3\mathbf{p}_0.
        \label{eq:faz_hacmi_korunumu}
    \end{equation}
\end{theorem}

\begin{corollary}[Geometric Origin of the Heisenberg Uncertainty Principle]
    \label{cor:heisenberg_kokeni_sc}
    The Heisenberg uncertainty relation:
    \begin{equation}
        \Delta x(t) \cdot \Delta p(t) = \left( \lambda(t) \Delta x_0 \right) \left( \frac{\Delta p_0}{\lambda(t)} \right) = \Delta x_0 \Delta p_0 \ge \frac{\hbar}{2},
        \label{eq:heisenberg_turetilis}
    \end{equation}
    is a direct consequence of the \textbf{topological conservation of Scalor phase space volume under conformal contraction}.
\end{corollary}

\section{Scale-Covariant Hamilton--Lagrange Mechanics and Dual Super-Manifolds}
\label{sec:lagrange_mekanigi}

Dynamical evolution and conservation laws in Scale-Time Calculus are derived variationally from the principle of least action ($\delta \mathcal{S} = 0$).

\subsection{Dual Super-Manifold $\mathcal{M}^{(4|1)}$ and the Berezin Integral}
To formally embed the nilpotent generator $\kappa$ ($\kappa^2 = 0$, Clifford~\cite{clifford1873} dual numbers) into differential geometry, we define a 5-dimensional **Dual Super-Manifold** $\mathcal{M}^{(4|1)}$ extending 4-dimensional Euclidean spacetime by a nilpotent Grassmann axis:
\begin{equation}
    Z^A \equiv \left( x^\mu, \theta \right), \quad \theta \equiv \kappa, \quad \kappa^2 = 0, \quad \mu \in \{0, 1, 2, 3\}.
    \label{eq:super_koordinat_sc}
\end{equation}
A dynamic Scalor super-field defined on this manifold takes the form:
\begin{equation}
    \mathbb{S}\left(x^\mu, \kappa\right) = \mathbf{X}_0^\mu(x) + \kappa \cdot \lambda(x).
    \label{eq:skalor_super_alan_sc}
\end{equation}
Using Berezin~\cite{berezin1966} integration rules:
\begin{equation}
    \int d\kappa \, 1 = 0, \qquad \int d\kappa \, \kappa = 1,
    \label{eq:berezin_kurallari_sc}
\end{equation}
the total Action Functional on the Dual Super-Manifold is defined as:
\begin{equation}
    \mathcal{S}[\mathbb{S}] = \int d^4x \int d\kappa \, \mathcal{L}\left(\mathbb{S}, \mathbb{D}_\mu \mathbb{S}\right).
    \label{eq:skalor_eylem_fonksiyoneli_sc}
\end{equation}

\subsection{Scale-Covariant Euler--Lagrange Equations}
Applying $\delta \mathcal{S} = 0$ yields the generalized \textbf{Scale-Covariant Euler--Lagrange Equations}:
\begin{equation}
    \mathbb{D}_\mu \left( \frac{\partial \mathcal{L}}{\partial (\mathbb{D}_\mu \mathbb{S})} \right) - \frac{\partial \mathcal{L}}{\partial \mathbb{S}} = 0.
    \label{eq:euler_lagrange_olcek}
\end{equation}
For a free Scalor field with Lagrangian density $\mathcal{L} = -\frac{1}{2} (\mathbb{D}_\mu \mathbb{S})(\mathbb{D}^\mu \mathbb{S}) - \frac{1}{2} m^2 \mathbb{S}^2$, Eq.~(\ref{eq:euler_lagrange_olcek}) analytically yields the Scale-Covariant Wave Equation~(\ref{eq:acik_dalga}).

\subsection{Scale-Covariant Noether Theorem and Energy-Momentum Conservation}
\label{sec:noether_ve_enerji}

In standard field theory, Noether's Theorem~\cite{noether1918} relates spacetime translation symmetries to conservation laws ($\partial_\mu T^{\mu\nu} = 0$). In a contracting universe where the atomic scale factor is time-dependent ($\lambda(t) = e^{-H_0 t}$), pure coordinate time translation ($t \to t + \delta t$) is not an isolated symmetry. Scale-Time Calculus resolves this directly by coupling scale and time:

\begin{definition}[Scale-Covariant Energy-Momentum Tensor]
    \label{def:olcek_enerji_momentum_tensoru_sc}
    The canonical energy-momentum tensor derived from diffeomorphism variations of the scale-covariant action is:
    \begin{equation}
        T^\mu_{\phantom{\mu}\nu} \equiv \frac{\partial \mathcal{L}}{\partial (\mathbb{D}_\mu \mathbb{S})} \mathbb{D}_\nu \mathbb{S} - \delta^\mu_\nu \mathcal{L}.
        \label{eq:olcek_enerji_momentum_tensoru}
    \end{equation}
\end{definition}

\begin{theorem}[Covariant Conservation of Energy-Momentum]
    \label{thm:kovaryant_korunum_sc}
    In an isolated contracting Scalor system, the Scale-Covariant divergence of the energy-momentum tensor vanishes identically:
    \begin{equation}
        \mathbb{D}_\mu T^{\mu\nu} = 0.
        \label{eq:kovaryant_korunum}
    \end{equation}
    Expanding in terms of partial derivatives on the flat background:
    \begin{equation}
        \partial_\mu T^{\mu\nu} + \kappa^\mu H_0 T_{\mu}^{\phantom{\mu}\nu} = 0.
    \end{equation}
    For the temporal component ($\nu = 0$), this yields the differential balance equation:
    \begin{equation}
        \frac{\partial T^{00}}{\partial t} + \nabla \cdot \mathbf{T}^{0} = - H_0 \left( T^{00} + \sum_{i=1}^3 T^{ii} \right).
        \label{eq:enerji_bilancosu}
    \end{equation}
\end{theorem}

\begin{proof}
    Consider the variation of the action under combined \textit{Scale-Time Translations}. Applying an infinitesimal coordinate shift $x^\mu \to x^\mu + \epsilon^\mu$ accompanied by the scale transformation $\lambda \to \lambda - \epsilon^\mu \kappa_\mu H_0 \lambda$, the field perturbation transforms covariantly as $\delta (\mathbb{D}_\mu \mathbb{S}) = \epsilon^\nu \mathbb{D}_\nu (\mathbb{D}_\mu \mathbb{S})$. Expanding via the chain rule and the Euler--Lagrange equations~(\ref{eq:euler_lagrange_olcek}):
    \begin{align*}
        \mathbb{D}_\mu \mathcal{L} & = \frac{\partial \mathcal{L}}{\partial \mathbb{S}} \mathbb{D}_\mu \mathbb{S} + \frac{\partial \mathcal{L}}{\partial (\mathbb{D}_\nu \mathbb{S})} \mathbb{D}_\mu (\mathbb{D}_\nu \mathbb{S})                                                               \\
                                    & = \left[ \mathbb{D}_\nu \left( \frac{\partial \mathcal{L}}{\partial (\mathbb{D}_\nu \mathbb{S})} \right) \right] \mathbb{D}_\mu \mathbb{S} + \frac{\partial \mathcal{L}}{\partial (\mathbb{D}_\nu \mathbb{S})} \mathbb{D}_\nu (\mathbb{D}_\mu \mathbb{S}) \\
                                    & = \mathbb{D}_\nu \left[ \frac{\partial \mathcal{L}}{\partial (\mathbb{D}_\nu \mathbb{S})} \mathbb{D}_\mu \mathbb{S} \right].
    \end{align*}
    Equating this with $\delta^\nu_\mu \mathbb{D}_\nu \mathcal{L}$ directly gives $\mathbb{D}_\nu T^\nu_{\phantom{\nu}\mu} = 0$. Due to nilpotency ($\kappa^2 = 0$), higher-order scale perturbations vanish identically, establishing exact covariant conservation.
\end{proof}

\begin{corollary}[Matter Mass Growth and Vacuum Flux Energy Balance]
    \label{cor:madde_vakum_enerji_dengesi_sc}
    Equation~(\ref{eq:enerji_bilancosu}) elucidates the physical exchange between matter contraction and the surrounding vacuum flow:
    \begin{enumerate}
        \item \textbf{Local Rest-Mass Energy Increase:} As $\lambda(t) \to 0$, conformal mass scaling dictates that particle rest energy increases as $E_{\mathrm{matter}}(t) = E_0 e^{H_0 t} \implies \frac{d E_{\mathrm{matter}}}{dt} = + H_0 E_{\mathrm{matter}}$.
        \item \textbf{Vacuum Flux Energy Supply:} The $-H_0 T^{00}$ term in Eq.~(\ref{eq:enerji_bilancosu}) demonstrates that this energy does not appear ex nihilo; it is precisely supplied by the radial vacuum inflow ($\mathcal{S}_{\mathrm{sink}}$) drawn from the surrounding continuum: $\frac{d E_{\mathrm{vacuum}}}{dt} = - H_0 E_{\mathrm{matter}}$.
        \item \textbf{Total Hyper-Energy Conservation:} When matter and its vacuum inflow basin are treated as a unified Scalor system, total energy is strictly conserved:
              \begin{equation}
                  \frac{d}{dt} \left( E_{\mathrm{matter}} + E_{\mathrm{vacuum}} \right) = 0 \iff \mathbb{D}_\mu T^{\mu\nu} = 0.
              \end{equation}
    \end{enumerate}
\end{corollary}

\section{Scale-Covariant Newtonian Mechanics and Orbital Stability}
\label{sec:newton_mekanigi}

The physical position vector defined with respect to the contracting atomic ruler is $\mathbf{R}(t) \equiv \mathbf{X}_0 / \lambda(t)$. Successive time derivatives yield:
\begin{align}
    \frac{d\mathbf{R}}{dt}     & = \frac{\dot{\mathbf{X}}_0}{\lambda} - \frac{\mathbf{X}_0 \dot{\lambda}}{\lambda^2} = \mathbf{v}_{\mathrm{local}} + H_0 \mathbf{R}, \label{eq:olcek_hiz} \\
    \frac{d^2\mathbf{R}}{dt^2} & = \mathbf{a}_{\mathrm{local}} + 2 H_0 \mathbf{v}_{\mathrm{local}} + H_0^2 \mathbf{R}. \label{eq:olcek_ivme}
\end{align}
The term $H_0 \mathbf{R}$ represents the kinematic scale recession drift arising from ruler contraction.

\begin{theorem}[Scale-Covariant Newton Equation of Motion]
    \label{thm:olcek_newton_hareket_sc}
    For a gravitational acceleration field $\mathbf{a}_{\mathrm{local}} = -\nabla V_{\mathrm{vacuum}}(\mathbf{R})$, the equation of motion in the contracting frame is:
    \begin{equation}
        m \frac{d^2\mathbf{R}}{dt^2} = - m \nabla V_{\mathrm{vacuum}}(\mathbf{R}) + 2 m H_0 \mathbf{v}_{\mathrm{local}} + m H_0^2 \mathbf{R}.
        \label{eq:olcek_newton_ana}
    \end{equation}
    This encompasses three physical contributions:
    \begin{enumerate}
        \item \textbf{Vacuum Inflow Force (Local Gravity):} $-m \nabla V_{\mathrm{vacuum}}(\mathbf{R})$.
        \item \textbf{Scale Damping Term:} $2 m H_0 \mathbf{v}_{\mathrm{local}}$ (cosmic Coriolis drag).
        \item \textbf{Apparent Cosmic Drift:} $m H_0^2 \mathbf{R}$ (corresponding to the geometric counterpart of Dark Energy / $\Lambda$).
    \end{enumerate}
\end{theorem}

\begin{corollary}[Local Locking vs. Cosmic Recession: Why the Moon Remains Bound]
    \label{cor:yerel_kilitlenme_kozmik_kaciklik_sc}
    Examining the distance scaling in Eq.~(\ref{eq:olcek_newton_ana}):
    \begin{itemize}
        \item \textbf{At Close Distances (Atomic to Lunar Scales):} Gravitational acceleration scales as $g_{\mathrm{vacuum}} = \frac{GM}{R^2} \propto R^{-2}$, whereas cosmic drift scales as $H_0^2 R \propto R$. At the Earth--Moon distance ($R \approx 3.84 \times 10^8\text{ m}$):
              \begin{equation}
                  g_{\mathrm{vacuum}} \approx 2.7 \times 10^{-3}\text{ m/s}^2 \quad \gg \quad H_0^2 R \approx 2.0 \times 10^{-27}\text{ m/s}^2.
              \end{equation}
              Gravitational binding dominates cosmic drift by \textbf{24 orders of magnitude ($10^{24}$ times)}. Consequently, the Moon is locked into its bound orbit, dynamically decoupling from cosmological expansion in accordance with the McVittie theorem. Solar System ephemeris telemetry limits are preserved.
        \item \textbf{At Intergalactic Distances ($R \gg 1\text{ Mpc}$):} Gravitational inflow decays to zero ($g_{\mathrm{vacuum}} \to 0$). Local locking vanishes, leaving only the Kinematic Scale Drift:
              \begin{equation}
                  \frac{d\mathbf{R}}{dt} \approx H_0 \mathbf{R} \quad \text{(Apparent Hubble Expansion)}.
              \end{equation}
    \end{itemize}
\end{corollary}

\section{Conclusion}
\label{sec:sonuc}

The Scale-Time Algebra formulated in this paper (the generator $\kappa$, Scalor space, and the operator $\mathbb{D}_t$) provides the requisite mathematical machinery to describe contracting matter in a static, non-expanding universe. Hubble's law is transformed from an empirical postulate into an exact calculus theorem, while photon propagation and atomic contraction are reconciled within a unified scale-covariant wave equation.

\begin{thebibliography}{99}
    \bibitem{main_paper}
    M.~F.~Haydari, ``The Unified Contraction Hypothesis and Scale-Time Calculus: A Conformal Geometric Treatise on Cosmology and Fundamental Interactions'', \textit{Zenodo Monograph}, 2026.

    \bibitem{paper1}
    M.~F.~Haydari, ``The Contracting Atom Metric and Conformal Action Principle: Testing Type Ia Supernovae (Pantheon+) and DESI 2024 BAO Without Dark Energy'' (Paper 1 in this series), \textit{Zenodo Preprints}, 2026.

    \bibitem{paper4}
    M.~F.~Haydari, ``Derivation of Gravitation as Spherical Vacuum Inflow Geometry: 3+1 ADM Action and the Painlevé--Gullstrand Metric'' (Paper 4 in this series), \textit{Zenodo Preprints}, 2026.

    \bibitem{paper5}
    M.~F.~Haydari, ``Electrodynamics, Vacuum Phase Waves, and Atomic Stability: Kinematic Pitch Interpretation of the Sommerfeld Constant $\alpha \approx 1/137$'' (Paper 5 in this series), \textit{Zenodo Preprints}, 2026.

    \bibitem{paper6}
    M.~F.~Haydari, ``Strong Nuclear Force, Quark Confinement, and Flux Tube Geometry in Scale-Time Calculus'' (Paper 6 in this series), \textit{Zenodo Preprints}, 2026.

    \bibitem{paper7}
    M.~F.~Haydari, ``Weak Nuclear Interaction, Chiral Phase Leakage, and Intermediate Vector Boson Scales'' (Paper 7 in this series), \textit{Zenodo Preprints}, 2026.

    \bibitem{wetterich2013}
    C.~Wetterich, ``A universe without expansion,'' \textit{Phys. Dark Univ.}, vol.~2, pp.~184--187, 2013.

    \bibitem{hoyle1964}
    F.~Hoyle and J.~V.~Narlikar, ``A new theory of gravitation,'' \textit{Proc. R. Soc. Lond. A}, vol.~282, no.~1389, pp.~191--212, 1964.

    \bibitem{weyl1918}
    H.~Weyl, ``Gravitation und Elektrizit\"at,'' \textit{Sitzungsber. Preuss. Akad. Wiss.}, pp.~465--480, 1918.

    \bibitem{clifford1873}
    W.~K.~Clifford, ``Preliminary sketch of biquaternions,'' \textit{Proc. London Math. Soc.}, vol.~4, pp.~381--395, 1873.

    \bibitem{berezin1966}
    F.~A.~Berezin, \textit{The Method of Second Quantization}, Academic Press, New York, 1966.

    \bibitem{noether1918}
    E.~Noether, ``Invariante Variationsprobleme,'' \textit{Nachr. Ges. Wiss. G\"ottingen, Math.-Phys. Kl.}, vol.~1918, pp.~235--257, 1918.
\end{thebibliography}

\end{document}
